\section{Introduction}\label{sec:current-overview}\label{sec:introduction}
A prepared experiment assigns a terminal probability law to each finite
word of commands. Even when its protected state has finite dimension,
an exact classical realization of all these laws can require infinitely
many retained labels. When a finite realization does exist, three sizes
must still be distinguished: the linear response dimension, the number
of distinct deterministic residuals, and the smallest stochastic memory.
The last can be strictly between the first two. We determine the form of
that stochastic advantage for reversible observable dynamics.

Let $b:A^*\to\Delta(C)$ be a probability response on finite command and
terminal alphabets. Its minimal linear space is spanned by the residual
functions $b^w(z)=b(wz)$. We call $b$ compactly reversible when its command
shifts are invertible on this space and generate a group with compact
closure. Prepared finite-dimensional unitary commands always have this
property after directions invisible to every future commanded terminal
test have been removed. The condition concerns observable responses;
the full physical unitary group may be larger.

The finite realization criterion in Theorem~\ref{thm:r72-reversible}
states that this compact observable group must be finite. The main new
result is more precise. Theorem~\ref{thm:r73-normal-form} replaces any
$M$-label stochastic realization by a permutation core and at most one
transient initializer, using at most $M$ labels in total. Its proof does
not merely apply a matrix-group theorem to the supplied transition
matrices: those matrices need not form a group, and their inaccessible
hidden responses need not belong to the minimal observable space. We
average words acting identically on the observable response, pass to the
closed classes of that finite chain, and prove that the compressed
command matrices factor through the observable group. The prescribed
initial state determines whether a further initializer is necessary.

The resulting variational formula, Theorem~\ref{thm:r73-memory-formula},
minimizes over finite permutation actions and bounded probability
decoders. It is an exact optimization over all real stochastic
competitors. In particular $|G_b|\le M!$, and a strict stochastic saving
over the deterministic residual count requires
$|G_b|\le(M-1)!$. Both bounds are sharp. A two-command family on $k$
coordinates has exactly $k+1$ stochastic live labels and $k!$
deterministic-transition labels for every $k\ge3$. The family has
rational preparations and readouts and admits a diagonal physical
realization. It exhibits a stochastic implementation advantage; it is
not presented as a quantum advantage over a classical experiment.

Probability constraints in the factorization carry information beyond the
group and its linear representation. Proposition~\ref{prop:r73-decoder-obstruction}
gives two responses with the same observable group $S_3$, linear dimension
$3$ and six deterministic residuals, but with stochastic minima $4$ and $6$.
The latter has an explicit protected-qubit realization. Its endpoint laws
force incompatible decoder supports in every action with at most five
points. Thus neither residual counting nor permutation-module dimension
alone determines the stochastic resource.

Theorem~\ref{thm:r73-algebraic-synthesis} makes the classification
effective for supplied algebraic observable data. A finite-group order
bound gives a terminating test of existence of any finite realization.
The normal form then reduces exact minimal synthesis to finitely many
finite-group actions and bounded bilinear probability equations.
The returned controller has algebraic entries and a counted initial
label. Theorem~\ref{thm:r73-finite-tests} and the ensuing gap theorem
give finite word certificates against arbitrary bounded-size stochastic
competitors. Applied to damping toward the group-averaged response, they
upgrade qualitative compactness to a computable positive interval on
which the exact minimal memory is unchanged. This interval is not
claimed to have the sharp first transition point.

\subsection*{Classical structure and the present assertion}
Finite residual equivalence for deterministic output machines is
classical \cite{RabinScott59}. Positive realization, observable operator
models and finite linear representations also have an extensive history
\cite{Heller65,CarlylePaz71,Jaeger00,ThonJaeger15,MonrasWinter14}.
Schwarz classified stochastic matrix groups, and Flor gave an elementary
description of nonnegative matrix groups, including the bounded-group
permutation structure and Brown's finiteness theorem
\cite{Schwarz64,Flor69,Brown64}. These results are prior mathematics.
Group images of completely simple semigroups also have a classical
structure theory \cite{Schwarz62}.
The issue addressed here is how an arbitrary stochastic realization of a
reversible \emph{observable} response can be compressed to such a group
while preserving a pure initial state, its terminal law, and the exact
label budget. The initializer-sensitive formula and its constructive
consequences are the statements for which a contribution is asserted.

Finite matrix-group recognition, the congruence bounds of Minkowski, and
real-algebraic automaton minimization remain credited to their existing
theories \cite{Serre07,BabaiBealsRockmore93,KieferWachter14}. The
finite-word equivalence argument is included as a classical lemma
\cite{Tzeng92}. Likewise, finite-dimensional quantum processes without
finite classical realizations and noise-dependent memory lower bounds
predate this paper \cite{Fanizza24}. Neither a general quantum/classical
separation nor those algorithmic ingredients are claimed as discoveries
of the present work. The statements and proofs below permit a direct
comparison with that literature; an exhaustive priority or editorial
verdict is not a mathematical consequence of them.

\subsection*{Relation to the causal resource theory}
The exact periodic partition law, its one-label leakage increment, and
the single-frequency and noncommuting-pair arithmetic envelopes remain
complete current results. Their proofs are in the paired supplement
except for the all-word reversible criterion and common-command
construction retained in the main article. The new normal form places
finite cyclic cases within a general finite-group realization problem,
while strictly damped irrational modes remain outside the reversible
class and retain their distinct spectral arguments.

The main article also retains the full actual-history causal comparison
and common paid construction. Its lower witness is formed on the
competing experiment's actual pre-call histories; it does not assume a
terminating reference run on another experiment. The supplement retains
performed-law compactness, boundary primal and dual attainment, robust
tied optimality, normalized stopping rigidity, instrument realization,
finite-width transfer, reset control, rank-constrained discrimination,
spectral geometry and composition. The independently compiled companion
retains absolute calibration and nonintervening directional sensing.
These results keep their individual hypotheses. Exact response
simulation is a component of this resource theory and is not a substitute
for a general classification of optimal feedback experiments.
